% FILENAME: paper/section-Conclusion.tex

\todo{Author Overview: The victory lap. Summarize that we didn't add anything; we took a pure gauge connection and derived gravity, matter, and QM simply by applying strict geometry.}

\todo{Author: Review and rename section title if desired.}
\section{Conclusion and Future Outlook}
\label{sec:Conclusion}

\todo{Author: Recap the core thesis. We started with zero free parameters. From a pure, background-independent gauge connection, we derived the Big Bang, the arrow of time, Ricci-flat gravity, dark matter, positive mass gaps, and the Schrödinger equation as absolute mathematical necessities.}

\todo{Author: The Paradigm Shift. Physics does not need extra dimensions or fundamental point particles. We do not need to "quantize gravity" because quantum mechanics is simply the relational phase space of gravity's topological defects measuring each other.}

\todo{Author: Future Outlook. Where does CGD go next? 
1. Computational: Feeding these exact geometric bounding rules into supercomputers to derive scattering cross-sections natively, without Feynman diagrams. 
2. Phenomenological: Extending the formal verification to the electroweak sector and the remaining Standard Model generations. 
3. Mathematical: Expanding the Litlib library to encompass the entirety of 20th-century differential geometry.}
